The accuracy of numerical methods for solving ordinary differential equations (ODEs) is typically analyzed by comparing the local truncation error, which measures the error introduced in a single step, and the global error, which accumulates over multiple steps. The fundamental approach involves Taylor series expansion and comparison with the exact solution.

Consider a general first-order ODE:
\begin{equation}
    \frac{dy}{dt} = f(t, y), \quad y(t_0) = y_0
\end{equation}

Let $y(t)$ be the exact solution and $y_n$ be the numerical approximation at time $t_n$. The local truncation error $\tau_n$ is defined as the error introduced in one step assuming previous steps were exact:
\begin{equation}
    \tau_n = y(t_{n+1}) - y_{n+1} \quad \text{given that} \quad y_n = y(t_n)
\end{equation}

The exact solution can be expanded around $t_n$:
\begin{equation}
    y(t_{n+1}) = y(t_n) + h y'(t_n) + \frac{h^2}{2} y''(t_n) + \frac{h^3}{6} y'''(t_n) + \mathcal{O}(h^4)
\end{equation}
where $h = t_{n+1} - t_n$ is the step size.

For a numerical method, we express the update formula as:
\begin{equation}
    y_{n+1} = y_n + h \Phi(t_n, y_n, h)
\end{equation}
where $\Phi$ represents the increment function specific to each method.

The local truncation error becomes:
\begin{equation}
    \tau_n = \left[y(t_n) + h y'(t_n) + \frac{h^2}{2} y''(t_n) + \cdots\right] - \left[y(t_n) + h \Phi(t_n, y(t_n), h)\right]
\end{equation}

A method is said to be of order $p$ if the local truncation error satisfies:
\begin{equation}
    \tau_n = \mathcal{O}(h^{p+1})
\end{equation}
which implies that the global error over a fixed time interval is $\mathcal{O}(h^p)$.

To determine the order $p$, we:
\begin{enumerate}
    \item Expand both the exact solution and numerical method using Taylor series
    \item Compare the expansions term by term
    \item Find the highest power of $h$ for which they agree
    \item The first non-matching term determines the local error order
\end{enumerate}

For the Hodgkin-Huxley system:
\begin{equation}
    C_m \frac{dV}{dt} = I_{ext} - g_{Na}m^3h(V-E_{Na}) - g_Kn^4(V-E_K) - g_L(V-E_L)
\end{equation}
with additional gating variable equations:
\begin{align}
    \frac{dm}{dt} &= \alpha_m(V)(1-m) - \beta_m(V)m \\
    \frac{dh}{dt} &= \alpha_h(V)(1-h) - \beta_h(V)h \\
    \frac{dn}{dt} &= \alpha_n(V)(1-n) - \beta_n(V)n
\end{align}

The accuracy analysis must consider:
\begin{itemize}
    \item The coupled nature of the four differential equations
    \item The nonlinear voltage dependence of rate constants $\alpha_x(V)$, $\beta_x(V)$
    \item The multiplicative terms ($m^3h$, $n^4$) in the current expressions
    \item The stiffness introduced by different time scales of gating variables
\end{itemize}

The local truncation error analysis follows the same Taylor expansion approach, but must account for these complexities through careful expansion of the nonlinear terms and their derivatives.
